\chapter{Conclusions}
\label{conclusions}

This thesis describes GJournal, a privacy-first iOS microjournaling app that runs all NLP on the device. Using the iOS Neural Processing Unit for on-device transformer inference proved viable across both embedding and emotion classification tasks.

The architecture combines quantized neural models, local SwiftData storage, and incremental clustering without sending any user data off device. This matters for a journaling app: users are unlikely to write honestly if their entries are processed by a remote server.

The experiments confirm that transformers can be fine-tuned, quantized, and deployed on a consumer iPhone without major accuracy loss on short-form text. Clustering also surfaces patterns from journal entries that can be browsed as calendar heatmaps, filtered lists, and a knowledge graph.

\par
The evaluation in Chapter~\ref{chap:evaluation} produced concrete numbers across all four models. all-MiniLM-L6-v2 achieved Spearman $r_s = 0.867$ on STS-B and V-measure 0.662 on the journal-domain validation set, outperforming Nomic-256 (V-measure 0.511), Nomic-768 (V-measure 0.481), and Apple NLEmbedding (V-measure 0.459) on clustering quality. Apple NLEmbedding scored Spearman 0.717 on STS-B, lower than all dedicated models but requiring zero additional bundle size. The MRL dimension trade-off analysis showed that Nomic inference latency stays flat at approximately 20\,ms across all truncation levels, while the Spearman drop from 768 to 256 dimensions is only 0.0044. The MobileBERT emotion classifier reached 59.5\% accuracy on a 12-class GoEmotions test set, with strong performance on the most common journaling emotions (Gratitude F1 0.813, Love F1 0.663, Anger F1 0.616).

\par
The primary contributions of this thesis are: the design of a complete end-to-end NLP pipeline that runs entirely on-device; the implementation of quantized transformer architectures for mobile execution; an incremental clustering approach tailored for short personal entries; a privacy-centric application architecture without cloud components; an empirical evaluation of Matryoshka Representation Learning truncation on journal-domain clustering, showing that 256-dim Nomic embeddings retain over 99.5\% of Spearman correlation relative to full 768-dim while losing 15 percentage points of V-measure against MiniLM on domain data; and the first benchmark of Apple NLEmbedding against dedicated sentence embedding models on semantic similarity and clustering tasks, conducted by bridging the Swift SDK into the Python evaluation pipeline via a standalone Swift script.

\section{Future Work}

A few areas remain open for follow-up work.

\textbf{Finer emotion taxonomy.} The 12-class GoEmotions mapping collapses related emotions to make the classification task tractable on mobile hardware. A more granular taxonomy, or one calibrated specifically to journaling language, would improve the quality of emotion-based filtering and the cluster-naming heuristic.

\textbf{Clustering stability.} DBSCAN recomputes all clusters from scratch when a new entry arrives. For a corpus of several hundred entries this is acceptable, but it becomes slow and produces label churn as the user adds entries over weeks. Incremental DBSCAN or a hierarchical approach that reuses previous cluster assignments would reduce both latency and label instability.

\textbf{Multimodal input.} The current system processes only text. Short audio recordings or attached photos are common in mobile journaling workflows. Extending the embedding pipeline to fuse text with speech or image features would broaden the kinds of patterns the system can surface.

\textbf{On-device MRL search.} The MRL truncation flexibility studied here was applied only to clustering. The same property makes MRL embeddings attractive for fast approximate nearest-neighbour search: a user could search past entries using a 64-dimensional index for speed, then re-rank results using the full 256-dim vectors. This retrieval pattern would complement the existing clustering view without requiring a second model.

\textbf{Schema migration and App Store readiness.} The SwiftData schema currently wipes stored data on model changes. A versioned migration policy, a privacy policy URL, and a VoiceOver accessibility pass are the remaining steps before an App Store submission.